\documentclass[12pt, letterpaper]{article} 

\usepackage{amsmath, amssymb, amsthm, enumerate}

\title{MATH 189 - Data Analysis}
\author{Dani Nguyen}
\date{Winter 2025}

\begin{document}
\maketitle
\tableofcontents

\section{Statistical Learning}
\subsection{Terminology}
"Supervised Leanring"  means we have both predictors and responses. \\
"Predictors/input variables/features" are arranged in a vector.
\begin{align*}
    X = (X_1,X_2,...,X_n)
\end{align*}
"Output variables/responses" is a scalar where $y \in \mathbb{R}$ \\
We think of data as coming from a model.
\begin{align*}
    y = f(x) + \epsilon
\end{align*}
$\begin{cases}
    f(x) \text{ is fixed but unknown function} \\
    \epsilon \text{ is the random error term with mean zero, independent of X}
\end{cases}$ \\ \\
Question: Why estimate $f$? \\
Prediction: Given $X$, we want to predict $y$ with $\hat{y} = \hat{f}(x)$ \\ \\ 
Mean square error is 
\begin{align*}
    \mathbb{E}(y-\hat{y})^2 &= \mathbb{E}(f(x)+\epsilon - \hat{f}(x)) \\
    &= \mathbb{E}(f(x)-\hat{f}(x))^2 + 2\mathbb{E}((f(x) - \hat{f}(x))\epsilon) + \mathbb{E}(\epsilon^2)
\end{align*}
Examine the middle term in detail 
\begin{align*}
    2\mathbb{E}((f(x) - \hat{f}(x))\epsilon)
    = \mathbb{E}(f(x)-\hat{f}(x))\cdot\mathbb{E}(\epsilon)=0
\end{align*}
We know that 
\begin{align*}
    \mathbb{E}(y-\hat{y})^2 = (f(x)-\hat{f}(x))^2 + \mathbb{E}(\epsilon^2)
\end{align*}
Where the left term is reducible error and the right term is irreducible error. \\ \\ 
"Inference" - Sometimes we want to understand the association between $y$ and $X_1, ... ,X_p$ 
\begin{enumerate}
    \item Which predictors are associated with response?
    \item Is the association positive or negative?
    \item Is the relationship approximately linear?
\end{enumerate}

\subsection{Accuracy}
\underline{Training data} 
\begin{align*}
    (x_1, y_1),(x_2, y_2),...(x_n, y_n) \\ 
    \text{where } X_i = (X_{i1}, X_{i2}, ... , X_{ip})
\end{align*}
n = \# data points (e.g. $n=30$) \\
p = \# predictors (e.g. $p=3$) \\ \\
\noindent \underline{Mean Square Error (MSE)}
\begin{align*}
    \text{MSE} = \frac 1n \sum_{i=1}^{n} (y_i - \hat{f}(x_i))^2 \\
    \text{Defined on training data} \\ 
    \text{Test MSE} = \text{ Ave } (y_0 - \hat{f}(x_0))^2 \\ 
    \text{Defined on previously unseen data point } (x_0, y_0)
\end{align*}
\underline{Linear / affine model}
\begin{align*}
    f(x) = \beta_0 + \beta_1x_1 + ... + \beta_px_p
\end{align*}
"Fit/train the model" usually means run an optimization on your computer
that uses training data and outputs the parameters $\beta_0, \beta_1,...,\beta_p$ that make the 
MSE and hopefully the test MSE small. \\

\noindent \underline{Classification} \\ 
Data set $(x_1,y_1),(x_2,y_2,),..., (x_n,y_n)$ where $X_i=(x_{i1},...,x_{ip})$ and $y_i$ is qualitative. \\ \\ 
Error rate is the fraction of incorrect classifications. 
\begin{align*}
    \frac 1n \sum_{i=1}^{n} \mathbb{I} \{y_i \ne \hat{y}_i\}
\end{align*}
$\mathbb{I}\{ y_i \ne \hat{y}_i \} = \begin{cases}
    1, & y_i \ne \hat y_i \\
    0, & y_i = \hat y_i 
\end{cases}$ \\ \\
Test error rate 
\begin{align*}
    \text{Ave}(\mathbb{I} \{ y_0 \ne \hat y_0 \})
\end{align*}
Where $(x_0,y_0)$ is a previously unseen test data point. \\ \\
\underline{K-nearest neighbors (KNN)} \\ 
The KNN classifier is defined by letting $\eta_0$ = set of indices $i\in\{ 1,...,n \}$ for K nearest neighbors of $X_0$ among the input data. \\\\
$\hat y_0 = j$ where $j$ maximizes 
\begin{align*}
    \sum_{i\in \eta_0}^{} \mathbb{I} \{y_i=j\}
\end{align*}
which is the number of nearest neighbors (out of K) assigned label $j$.

\section{Linear Regression}
\subsection{Simple Linear Regression}
Underlying statistical model: $y=\beta_0 + \beta_1X + \epsilon$
where $y$ is the quantitative response and a single predictor $X$ which is why it is called simple. \\

\noindent Prediction model: $\hat y = \hat \beta_0 + \hat \beta_1X$ \\
$\hat \beta_0, \hat \beta_1$ are estimates based on training data, where hat means estimated quantity. \\

\noindent \underline{Estimating coefficients} \\
The residual sum of squares (RSS) 
\begin{align*}
    \text{RSS} &= \sum_{i=1}^{n}(y_i-\hat y_i)^2 \\
    &= \sum_{i=1}^{n}(y_i-\hat \beta_0 - \hat \beta_1X)^2 
\end{align*}
Choose $\hat \beta_0, \hat \beta_1$ to minimize RSS.
\begin{align*}
    \frac{\partial \text{RSS}}{\partial \hat\beta_0} &= \frac{\partial}{\partial \hat\beta_0}  \sum_{i=1}^{n}(y_i-\hat \beta_0 - \hat \beta_1X)^2  \\
    &= 2\sum_{i=1}^{n}(\hat \beta_0+\hat \beta_1X_i-Y_i) \\
    &\implies \text{RSS is } \begin{cases}
        \text{decreasing for really small } \hat \beta_0 \\
        \text{increasing for really big } \hat \beta_0 
    \end{cases} 
\end{align*}
RSS is optimal (in $\hat \beta_0$) when 
\begin{align*}
    0 &= \frac 1n \sum_{i=1}^{n} (\hat \beta_0 + \hat \beta_1 X_i - Y_i) \\
    &= \hat \beta_0 + \hat \beta_1 \Big(\frac 1n \sum_{i=1}^{n}X_i\Big) - \Big(\frac 1n \sum_{i=1}^{n}Y_i\Big) \\ 
    &= \hat \beta_0 +\hat \beta \bar x - \bar y \\
    \hat \beta_0 &=\bar y - \hat \beta_1 \bar x 
\end{align*}
Where $\begin{cases}
    \bar x = \frac 1n \sum_{i=1}^{n}X_i \text{ is the mean predictor} \\
    \bar y = \frac 1n \sum_{i=1}^{n}Y_i \text{ is the mean response}
\end{cases}$ \\  \\
\underline{Covariance}
\begin{align*}
    \text{Cov}  (X_i, Y_i)^n_{i=1} &= \frac 1n \sum_{i=1}^{n}X_iY_i - \bar X \bar Y \\
    &= \frac 1n \sum_{i=1}^{n} (X_i - \bar x) (Y_i - \bar y)
\end{align*}
\underline{Variance}
    \begin{align*}
        \text{Var}(X_i)^n_{i=1} &= \text{Cov} (X_i,X_i)_{i=1}^n \\
        &= \frac 1n \sum_{i=1}^{n}(X_i - \bar x)^2
    \end{align*}
From this we can get our optimal $\hat \beta_1$
\begin{align*}
    \hat \beta_1 &= \frac{\text{Cov}(X_i,Y_i)_{i=1}^n}{\text{Var}(X_i)_{i=1}^n} \\ 
    &= \frac{\frac 1n \sum_{i=1}^{n}(X_i-\bar x)(Y_i-\bar y)}{\frac 1n \sum_{i=1}^{n} (X_i-\bar x)^2}
\end{align*}

\end{document}